% Cumulative fragment to insert after the presentation of the
% TPK update in nucleo.tex. Does not replace nucleo, extension,
% generacion, or registro_k. Detailed provenance in the companion MD.
\subsection{Operator architecture of transport, memory, and prolongation}
\label{ii:tpk-arquitectura}

The topological phase kernel organizes the evolution of the two APP
evaluations on oriented histories. Its structure brings together a path support,
local update rules, prolongation relations, temporal
records, and reading maps. Each component has its own
domain. This hierarchy allows the observable position, the transport
quotient, and the memory of the crossed boundaries to be preserved within
a single construction. The composition in equation~\eqref{eq:actualizacion}
thus acquires a precise realization at each of these levels.

\subsubsection{Bivariate paths and lifted transport}

Let \(\mathsf P_{\rm APP}\) be the free category of the cardinal graph on
\(X_9=(\mathbb Z/9\mathbb Z)^2\). Its objects are the APP positions;
its morphisms are cardinal words, with composition by
concatenation. The category
\[
 \mathsf P_{\rm APP}^{(2)}
 =\mathsf P_{\rm APP}\times\mathsf P_{\rm APP}
\]
preserves, in order, the path of the additive sheet and that of the
multiplicative sheet. Both paths traverse the same arithmetic support and
keep their evaluations distinct.

In the chart \(I_9=\{0,\ldots,8\}\), the generators are
\[
 \nu_N=(-1,0),\quad\nu_E=(0,1),\quad
 \nu_S=(1,0),\quad\nu_O=(0,-1),
 \qquad T_d(x)=[x+\nu_d]_9.
\]
For \(w=d_1\cdots d_r\), the lift to the integer covering space preserves
\[
 \widetilde x_k=\widetilde x_0+\sum_{j=1}^{k}\nu_{d_j},
 \qquad x_k=[\widetilde x_k]_9.
\]
When \(x_r=x_0\), the winding class is
\begin{equation}
 \operatorname{wind}(w)
 =\frac{1}{9}(\#S-\#N,\#E-\#O)\in\mathbb Z^2.
 \label{ii:tpk-enrollamiento}
\end{equation}
The final position and winding are two different readings of the
same route. The second records the displacement of the lift even
when the first has returned.

Path reversal has a determined domain and codomain:
\[
 \operatorname{rev}:\operatorname{Hom}(x,y)
 \longrightarrow\operatorname{Hom}(y,x),\qquad
 d_1\cdots d_r\longmapsto\overline d_r\cdots\overline d_1,
\]
where \(\overline N=S\), \(\overline E=O\). It satisfies
\(\operatorname{rev}^2=I\), reverses the order of composition, and changes the
sign of~\eqref{ii:tpk-enrollamiento}. This is the precise operation that
must be transported when comparing conjugate orientations of a route.

\subsubsection{Full state and preservation through updates}

The description of \(\widehat X\) is organized over the positions and
directions of the two cursors. Over each observable configuration
\(q\), the arithmetic fiber is written
\begin{equation}
 \mathcal F_q=
 \mathsf O_q\times\mathsf H_q\times\mathsf C_q\times
 \mathsf S_q\times\mathsf B_q\times\symup{\Lambda}_q.
 \label{ii:tpk-fibra}
\end{equation}
The coordinates \(\mathsf O_q\) preserve orientation and sheet;
\(\mathsf H_q\), the remainders and quotients of the evaluations;
\(\mathsf C_q\), compatible carries;
\(\mathsf S_q\), the boundary signature;
\(\mathsf B_q\), the prefix and its cylinders; and
\(\symup{\Lambda}_q\), the ordered transport record. The fiber is
restricted by the corresponding arithmetic, incidence, and
truncation identities. The full state also preserves the route
that produced it and the origin from which its phases are read.

Mode selection has type
\[
 M_{\rm ph}:D_9\longrightarrow\{-1,0,+1\},\qquad
 \mathsf{Sel}_t:\widehat X\longrightarrow\widehat X^{\rm sel}_t.
\]
The second map incorporates the value of the first into the state and
activates the corresponding sheet. Lifted transport
\[
 \mathsf{Tra}_t:\widehat X^{\rm sel}_t
 \longrightarrow\widehat X^{\rm tra}_t
\]
applies \(T_d\) to the active cursor, concatenates its route letter, and preserves the
selected sheet. Finally,
\[
 \mathsf{Upd}_t:\widehat X^{\rm tra}_t\longrightarrow\widehat X
\]
updates the evaluations and records. In this notation,
\(\mathsf{Upd}_t\) is the update \(\mathsf{Mem}_t\) of
\eqref{eq:actualizacion}, made explicit in its components.

Each integer evaluation \(a\) is preserved through
\(a=\rho_9(a)+9q_9(a)\); each ternary evaluation, through its oriented
remainder and carry. For a concatenation \(r_2r_1\), carry
transport satisfies the cocycle law
\begin{equation}
 k_{\rm tr}(r_2r_1)
 =k_{\rm tr}(r_2)+\mathsf C(r_2)k_{\rm tr}(r_1),
 \label{ii:tpk-cociclo-transporte}
\end{equation}
where \(\mathsf C(r_2)\) is the transport action of the second route
on the carry coordinate of the first. In the ordered record,
updating is performed by adjoining the event produced:
\[
 \Lambda_{t+1}=\Lambda_t\Vert
 (\text{sheet},d_t,\text{remainder},\text{quotient},
   \text{carry},\text{orientation},\text{boundary})_t.
\]
Preservation has a dynamic meaning here: earlier contributions
remain recoverable within the record being prolonged.
Their coordinates may change when the state is transported.

\subsubsection{Affine action on transport memory}

For a trajectory \(r:q\to q'\), the maps
\[
 \mathsf H(r):\mathsf H_q\to\mathsf H_{q'},\qquad
 \mathsf C(r):\mathsf C_q\to\mathsf C_{q'}
\]
respect identity and concatenation; \(\mathsf C_q\) is the abelian
group of transport memory, and \(\mathsf C(r)\) its homomorphism.
The increment \(k_{\rm tr}(r)\in\mathsf C_{q'}\), with
\(k_{\rm tr}(\operatorname{id}_q)=0\), is the cocycle of
\eqref{ii:tpk-cociclo-transporte}. This symbol distinguishes the transport
increment from the rational coordinates of the dodecaphasic record.

In a chart with data \((h,c,\lambda)\), the update has the
explicit form
\begin{equation}
 h'=\mathsf H(r)h,\qquad
 c'=\mathsf C(r)c+k_{\rm tr}(r),\qquad
 \lambda'=r\circ\lambda.
 \label{ii:tpk-accion-afin}
\end{equation}
The signature is evaluated through
\(\operatorname{Sig}_{q'}(h',c',\lambda')\). Its residence in the
fiber \(\mathsf S_{q'}\) preserves its dependence on those data.

\begin{proposition}[Composition of affine memory]
\label{ii:tpk-composicion-afin}
Updating by \(r_1\), followed by updating by
\(r_2\), agrees with updating by \(r_2r_1\).
\end{proposition}
\begin{proof}
Two memory updates give
\[
 c''=\mathsf C(r_2)\mathsf C(r_1)c
       +\mathsf C(r_2)k_{\rm tr}(r_1)+k_{\rm tr}(r_2).
\]
Compatibility of \(\mathsf C\) with concatenation and
identity~\eqref{ii:tpk-cociclo-transporte} turn this expression
into \(\mathsf C(r_2r_1)c+k_{\rm tr}(r_2r_1)\).
The coordinate \(h\) satisfies the same compatibility by
functoriality of \(\mathsf H\); the trajectory does so by
associativity. The signature agrees because it is evaluated on the same data.
\end{proof}

The spatial chart of each cursor provides an elementary realization
of this law. For \(x\in I_9^2\) and a cardinal direction \(d\),
\[
 a_d(x)=\frac{x+\nu_d-T_d(x)}9\in\mathbb Z^2,
 \qquad x'=T_d(x),\qquad c'=c+a_d(x).
\]
Here \(\mathsf C(r)=I\), and \(k_{\rm tr}(r)\) sums the
integer boundary crossings. For a route of length \(n\),
the telescoping sum gives
\[
 x_n=x_0+\sum_{j=1}^n\nu_{d_j}
             -9\sum_{j=1}^n a_{d_j}(x_{j-1}).
\]
On a loop, the last sum equals
\(\operatorname{wind}(r)\). The residual position returns while
the lift preserves the number and orientation of its crossings.
The remaining memory components retain their own transport
maps and prolongation relations.

\subsubsection{Chronology of the two cursors and observable return}

The finite calendar lives in \(\Theta_{108}=\mathbb Z/108\mathbb Z\).
For \(0\leq\widetilde\theta<108\), set
\[
 r(\theta)=1+(\widetilde\theta\bmod9),\qquad
 \beta(\theta)=\left[\left\lfloor\frac{\widetilde\theta}{9}
 \right\rfloor\right]_{12},\qquad
 \varepsilon_9(r)=M_{\rm ph}(r).
\]
Presentation~\eqref{tpk-int:qobs} groups position and direction by
cursor. This extension groups the two positions first and
then the two directions. The two charts are identified through the
factor permutation
\begin{equation}
 \begin{aligned}
 \mathfrak p_{\rm obs}:\ (X_9\times\mathcal D)^2\times\Theta_{108}
 &\longrightarrow X_9^2\times\mathcal D^2\times\Theta_{108},\\
 (x_+,d_+,x_\times,d_\times,\theta)
 &\longmapsto(x_+,x_\times,d_+,d_\times,\theta),
 \end{aligned}
 \label{ii:permutacion-coordenadas-observables}
\end{equation}
with \(\mathcal D=\{N,E,S,O\}\). The inverse restores each direction alongside
its cursor. If superscripts \({\rm agr}\) and \({\rm sep}\) distinguish
these two expressions of the observable successor, then
\[
 \mathfrak p_{\rm obs}F_{\rm obs}^{\rm agr}
 =F_{\rm obs}^{\rm sep}\mathfrak p_{\rm obs}.
\]
Indeed, the selector depends on the same phase, the cardinal vector is
applied to the same cursor, and direction reversal occurs at the same
instant; only the coordinate order changes. We retain
\(\mathcal Q_{\rm obs}\) and \(F_{\rm obs}\) for the space and
map so identified. This presentation permutation leaves
the events and memory coordinates of the history unchanged.

The observable space is
\[
 \mathcal Q_{\rm obs}
 =X_9^2\times\{N,E,S,O\}^2\times\Theta_{108}.
\]
Write a state as
\(q=(x_+,x_\times,d_+,d_\times,\theta)\). The active sheet is first evaluated
at the current position. The positions are then updated by
\begin{equation}
 (x_+',x_\times')=
 \begin{cases}
 (T_{d_+}x_+,x_\times),&\varepsilon_9(r)=+1,\\
 (x_+,T_{d_\times}x_\times),&\varepsilon_9(r)=-1,\\
 (x_+,x_\times),&\varepsilon_9(r)=0.
 \end{cases}
 \label{ii:tpk-fobs-posiciones}
\end{equation}
If \(\widetilde\theta+1\) is a multiple of \(27\),
\((d_+,d_\times)\) is replaced by \((\overline d_+,\overline d_\times)\).
In all cases \(\theta\) is incremented modulo \(108\). These rules
define the successor
\(F_{\rm obs}:\mathcal Q_{\rm obs}\to\mathcal Q_{\rm obs}\).
The emission realization in section~\ref{sec:extension} applies
exactly this order of evaluation, transport, and reversal.

\begin{proposition}[Observable return and preservation of history]
\label{ii:tpk-retorno-observable}
The map \(F_{\rm obs}\) is a permutation. The operator
\(\mathcal K_{27}=F_{\rm obs}^{27}\) satisfies
\(\mathcal K_{27}^{4}=I\). A complete history of one hundred eight
updates also preserves the events and carry accumulated
during that turn.
\end{proposition}
\begin{proof}
Knowing the later phase recovers the preceding phase. The rule indicates
whether a direction reversal occurred; undoing it and applying
inverse cardinal transport to the active cursor recovers the previous state.
Each segment of twenty-seven transitions contains nine activations of each
sheet. From an endpoint of a segment, each cursor makes nine displacements
in the same direction, whose sum is zero in \(X_9\); it then reverses
its orientation. Two segments restore the orientations, and four also restore
the calendar. The identity is transported to any phase by
invertibility of \(F_{\rm obs}\). The full record preserves the
concatenation of the one hundred eight events and their counter, which have been
updated at every transition.
\end{proof}

The periodicity of this permutation belongs to the observable quotient. In
the full history, the integer window counter increases by twelve
over the same interval. Exact sequence
\eqref{eq:extension108} expresses the reduction of the counter to its
dodecaphasic coordinate and preserves its defect of additivity.

\subsubsection{Bidirectional electronic transport and arithmetic record}
\label{ii:transporte-electronico}

The central representation uses a cell shared by the two
APP evaluations. Its coordinates belong to \(D_9=\{1,\ldots,9\}\),
and the inversion \((r,c)\mapsto(10-r,10-c)\) has the unique fixed point
\((5,5)\). The visible state is written \((r,c,h,v,n)\), with
\(h,v\in\{-1,1\}\) and integer chronological index \(n\geq0\).
The tritic selector repeats \((+1,-1,0)\). On the additive sheet,
\(r+c\) is read and the column is translated in direction \(h\); on the
multiplicative sheet, \(rc\) is read and the row is translated in direction
\(v\). The threshold preserves position. Both orientations are
reversed upon completing each interval of twenty-seven updates.

To express reading and carry with the same positive convention,
define
\[
 [z]^+_9=1+((z-1)\bmod9),\qquad
 q_9(z)=\frac{z-[z]^+_9}{9}.
\]
Evaluation precedes displacement. If \(a_n\) is the integer
read, the event preserves \((a_n,[a_n]^+_9,q_9(a_n))\), the sheet,
previous position, orientation, and boundary crossing. The
tritic symbol is \([a_n]^+_9\bmod3\). The spatial lift
retains, for example, the integer
\(q_9(c+h)\) when \(c\) is replaced by \([c+h]^+_9\).

The shared-cell representation and the two-cursor
representation of \eqref{ii:tpk-fobs-posiciones} have different visible
states. Their records are related to the TPK state through
their respective representation maps. For the central trajectory,
the full update on the shared cell producing the electronic word
is used here.

\begin{proposition}[Invertible update with preserved quotients]
\label{ii:electron-inversa-cocientes}
Add to the visible state two integer records \(Q_+,Q_\times\)
and two winding numbers \(W_r,W_c\). Each active reading
increments its sheet's record by \(q_9(a_n)\); each displacement
increments the corresponding winding by its boundary crossing.
This update has an explicit inverse on its image.
\end{proposition}
\begin{proof}
The later index determines the preceding mode and whether an orientation
reversal has just occurred. This reversal is undone first.
If the mode was additive, the previous column is recovered as
\([c-h]^+_9\); if multiplicative, the row is recovered as
\([r-v]^+_9\). The neutral mode preserves both coordinates. The
recovered cell again determines the integer read, its remainder,
quotient, and boundary crossing. Subtracting these increments from
\(Q_+,Q_\times,W_r,W_c\) and reversing the chronological index
recovers the preceding state exactly. Composing these
inverses recovers an entire finite history. Event concatenation
also preserves the order of readings.
\end{proof}

The records \(Q_+,Q_\times\) are an explicit representation by sums
of the quotient history. The construction is integrated into the memory
fiber alongside its other coordinates. Its composition law,
for concatenable histories \(u,v\), is
\[
 Q_s(uv)=Q_s(u)+Q_s(v),\qquad s\in\{+,\times\},
\]
where the second history begins at the first one's final state.
Algebraic inversion of an update subtracts its increments;
subsequently traversing a spatial return path records
new evaluations. This distinction simultaneously preserves
transport reversibility and history duration.

\begin{proposition}[Central return and exact accumulation]
\label{ii:electron-memoria-retorno}
Starting from \((r,c,h,v,n)=(5,5,1,1,0)\) and zero initial records, the
spatial projection and its orientations return after fifty-four
updates. Over that interval,
\[
 (\Delta W_r,\Delta W_c)=(0,0),\qquad
 (\Delta Q_+,\Delta Q_\times)=(10,46).
\]
After \(54N\) updates, for every integer \(N\geq0\),
one has \((Q_+,Q_\times)=(10N,46N)\), while the spatial
coordinates and their orientations return. The integer index takes the value
\(54N\); the dodecaphasic calendar modulo \(108\) returns for
even \(N\).
\end{proposition}
\begin{proof}
Each pair of active modes displaces each coordinate once. In
twenty-seven updates, both complete a turn of nine
positions and return to the central diagonal; \(h,v\) are then
reversed. The next block traverses the opposite turn. Spatial
crossings cancel between the two blocks.

In each block, the additive readings run through \(2k\),
\(1\leq k\leq9\), whose positive quotient sums to five. The
multiplicative readings of the positive block run through
\(k[k+1]^+_9\). Their integer sum and remainder sum are
\[
 \sum_{k=1}^9 k[k+1]^+_9=249,\qquad
 \sum_{k=1}^9[k[k+1]^+_9]^+_9=42.
\]
The second identity follows from the list
\((2,6,3,2,3,6,2,9,9)\). By Euclidean division,
the quotient sum is \((249-42)/9=23\). Substitution
\(k\mapsto[k-1]^+_9\) shows that the negative block contains
the same multiset of products. The two blocks therefore produce
ten and forty-six units in the respective records.
The spatial projection and its orientations return to their initial
values. The activation and reversal rules repeat upon
shifting the index by \(54\), and the increments do not depend on
the accumulated records. The preceding additive law
proves the formula for every \(N\) by induction.
\end{proof}

The tritic reading and window orientation provide, on
the same trajectory,
\[
 \gamma_e=\bigl((100000220)^3(120200000)^3\bigr)^2,\qquad
 \tau_e=(+,+,+,-,-,-,+,+,+,-,-,-).
\]
Powers represent concatenations. The discrepancy
and window readers developed in
\eqref{eq:el-kd} and \eqref{eq:el-komega} both compute
\((80,54,6)\). The evaluation is obtained from the trajectory before
incorporating \(\alpha\), the action section, or a mass.
Quotient memory preserves information additional to that triple:
in one hundred eight updates it records \((20,92)\), with
zero net spatial winding. Each representation thus retains the
type of the observable producing it.


% Separate development, 9 September 2026.
% Extends the product reader of centro_electronico_registros.tex.
% Does not identify the spatial half-turn with a physical conjugation.
\subsubsection{Bidirectionality of the product reader and quotient memory}
\label{bim09:seccion}

The product reading preserves four arithmetic coordinates before
determining a signature. For $x=(i,j)\in D_9^2$, $D_9=\{1,\ldots,9\}$,
define
\[
 N(x)=ij,\qquad r(x)=1+((ij-1)\bmod9),\qquad
 q(x)=\frac{N(x)-r(x)}9,
\]
and the discrete reader
\[
 \ell(1,2,4,8,7,5)=(0,1,2,3,4,5),\qquad
 \ell(3)=\ell(6)=\ell(9)=0.
\]
Thus $N=r+9q$ simultaneously preserves the integer product, positive residue,
and arithmetic quotient. The function $f_\ell(x)=\ell(r(x))$ is the
coordinate used by the regional signature. The functions $N,r,q$, and
$f_\ell$ have a positional domain; the transport direction remains
in the oriented seed.

\paragraph{Advance, half-turn, and reading before transport.}

Let $\mathscr X_\times=D_9^2\times\{N,E,S,O\}$. The cardinal vectors
are $\nu_N=(-1,0)$, $\nu_E=(0,1)$, $\nu_S=(1,0)$,
$\nu_O=(0,-1)$. Positional addition is reduced to positive representatives
modulo nine. Define
\[
 P(x,d)=(x+\nu_d,d),\qquad H(x,d)=(x,\bar d),
 \qquad \nu_{\bar d}=-\nu_d.
\]
These operators satisfy $P^9=I$, $H^2=I$, and $HPH=P^{-1}$.

In the first fifty-four TPK transitions, the product channel
is activated in phases $2,5,8$ of each nonadic window. Each activation
reads the position before advance. After transitions $27$ and $54$,
$H$ is applied. Thus the eighteen positions read are, for a
seed $s$,
\[
 P^0s,P^1s,\ldots,P^8s,\quad
 HP^9s,\,P H P^9s,\ldots,P^8H P^9s,
\]
with the understanding that the positional reading omits only direction.
For any positional function $f$, put $f_n=f(P^ns)$, with
indices modulo nine. The word of raw readings is
\begin{equation}
 \mathcal R_f(s)=
 (f_0,f_1,\ldots,f_8,\ f_0,f_8,f_7,\ldots,f_1).
 \label{bim09:palabra}
\end{equation}
The six window sums, each comprising three activations, are
\begin{align}
 A_j^f(s)&=f_{3j}+f_{3j+1}+f_{3j+2},&&0\leq j<3,\label{bim09:ida}\\
 A_{3+j}^f(s)&=f_{-3j}+f_{-3j-1}+f_{-3j-2},&&0\leq j<3.
 \label{bim09:vuelta}
\end{align}
For $f=f_\ell$, the signature is
$E_{\rm sig}(s)=([A_0^f]_{10},\ldots,[A_5^f]_{10})$.
The decimal quotient of each accumulator is
$c_j^{10}=\lfloor A_j^f/10\rfloor$, and it preserves
$A_j^f=[A_j^f]_{10}+10c_j^{10}$.

\begin{proposition}[Half-turn and permutation of half-cycles]
\label{bim09:media-vuelta}
For every oriented seed and every positional function $f$,
\[
 \mathcal R_f(Hs)=\operatorname{rot}_9\mathcal R_f(s),\qquad
 A_j^f(Hs)=A_{j+3\bmod6}^f(s).
\]
The same permutation transports the decimal residues and quotients
of the accumulators. It applies, in particular, to $N,r,q$, and $f_\ell$.
\end{proposition}
\begin{proof}
The relation $P^nH=HP^{-n}$ and equality $f(Hs)=f(s)$ give
$f(P^nHs)=f(P^{-n}s)=f_{-n}$. Substitution into
\eqref{bim09:palabra} exchanges its two groups of nine entries.
Summation in groups of three produces the window identity. Euclidean
division by ten is performed componentwise and respects the
permutation.
\end{proof}

Advance also has an exact law on the accumulators:
\begin{align}
 A_j^f(Ps)-A_j^f(s)&=f_{3j+3}-f_{3j},&&0\leq j<3,
 \label{bim09:avance-ida}\\
 A_{3+j}^f(Ps)-A_{3+j}^f(s)&=f_{1-3j}-f_{-2-3j},&&0\leq j<3.
 \label{bim09:avance-vuelta}
\end{align}
The proof consists in canceling the two common summands of each window.
Transport of a signature therefore requires the boundary evaluations
appearing in these differences, in addition to the sums already
expressed.

\paragraph{Chronological inversion and boundary terms.}

Let $\gamma=(x_0,\ldots,x_n)$ be a recorded trajectory, and let
$\gamma^\dagger=(x_n,\ldots,x_0)$ be its chronological inversion, with each
edge oriented in the opposite direction. For a positional function $f$,
the reading before each advance is
\[
 C_f(\gamma)=\sum_{t=0}^{n-1}f(x_t).
\]
Then
\begin{equation}
 C_f(\gamma^\dagger)=C_f(\gamma)+f(x_n)-f(x_0).
 \label{bim09:frontera-inversion}
\end{equation}
If $n=3m$ and $A_j^f$ groups three consecutive edges,
\begin{equation}
 A_j^f(\gamma^\dagger)=A_{m-1-j}^f(\gamma)
 +f(x_{3(m-j)})-f(x_{3(m-j-1)}).
 \label{bim09:frontera-ventana}
\end{equation}
Both formulas follow because the reverse trajectory reads the arrival
endpoints of the original edges. The interior summands coincide;
the difference preserves the endpoints.

For an antisymmetric edge observable $a(x,y)=-a(y,x)$, its word
does transform according to
\[
 \mathcal R_a(\gamma^\dagger)
 =-\operatorname{rev}\mathcal R_a(\gamma).
\]
This identity acts on oriented edges. Identities
\eqref{bim09:frontera-inversion}--\eqref{bim09:frontera-ventana}
act on positional readings. The half-turn $H$ changes the
initial direction and preserves calendar order; chronological
inversion reverses the complete record. Their domains and reading rules
are specified separately.

\paragraph{Occupation cocycle and memory preserved on outward and return paths.}

For a cardinal edge $x\to y$, with positive representatives
$\widetilde x,\widetilde y$, the covering crossing is
\[
 b(x,y)=\frac{\widetilde x+\nu_d-\widetilde y}{9}\in\mathbb Z^2.
\]
The identity $b(y,x)=-b(x,y)$ holds. In the free category of recorded
paths, define the additive reader
\begin{equation}
 \mathfrak m(\gamma)=
 \left(\sum_{t=0}^{n-1}b(x_t,x_{t+1}),\ C_q(\gamma),\ n\right).
 \label{bim09:memoria-aditiva}
\end{equation}
Concatenation satisfies
$\mathfrak m(\gamma_2\circ\gamma_1)=
\mathfrak m(\gamma_1)+\mathfrak m(\gamma_2)$.
It is an additive cocycle on recorded trajectories: the record preserves
outward and return edges as different events.

By \eqref{bim09:frontera-inversion},
\begin{equation}
 \mathfrak m(\gamma^\dagger\circ\gamma)=
 \bigl(0,\ 2C_q(\gamma)+q(x_n)-q(x_0),\ 2n\bigr),
 \label{bim09:memoria-retorno}
\end{equation}
where $\gamma^\dagger\circ\gamma$ denotes the temporal concatenation
of the outward path followed by the return. The oriented sum of crossings cancels;
arithmetic occupation and length remain. This reader does not
descend to the quotient that identifies a path followed by its inverse
with the empty trajectory, since that quotient eliminates the events
preserved by the second and third coordinates.

There is also a purely boundary-based even reading:
\[
 V_b(\gamma)=\sum_{t=0}^{n-1}
 b(x_t,x_{t+1})b(x_t,x_{t+1})^{\mathsf T},\qquad
 V_b(\gamma^\dagger\circ\gamma)=2V_b(\gamma).
\]
The odd and even readings record distinct properties of
the same edges; both are computed before the path is forgotten.

\paragraph{Exact application to the signature fiber \texorpdfstring{$555555$}{555555}.}

The finite enumeration of $\mathscr X_\times$ contains $324$ seeds,
$26$ signatures, and $24$ seeds with signature $555555$. The stable
partition under $P,H$ refines this fiber into three classes of eight elements.
They can be labeled by their signature after one advance:
\[
 \mathscr C_{177},\quad\mathscr C_{717},\quad\mathscr C_{771}.
\]
The half-turn action is
\[
 H\mathscr C_{177}=\mathscr C_{177},\qquad
 H\mathscr C_{717}=\mathscr C_{771},\qquad
 H\mathscr C_{771}=\mathscr C_{717}.
\]
The attached finite certificate verifies these identities on all
representatives. In the twenty-four seeds, the six raw accumulators
of $f_\ell$ are exactly five. Hence the decimal quotients
$c_j^{10}$ vanish on this fiber. The additional information
is observed in the positions, integer evaluations, and arithmetic
quotients, which have readings different from that signature.

The fixed transverse coordinate $a$ of the cardinal path is $4$ or
$5$ in these seeds. During one turn of nine advances, each
value of the moving coordinate $u\in D_9$ is traversed once. Since
$\gcd(a,9)=1$, the positive residues $r(au)$ permute $D_9$.
This gives, without using a physical constant,
\[
 \sum_{u=1}^9au=45a,\qquad
 \sum_{u=1}^9r(au)=45,\qquad
 \sum_{u=1}^9q(au)=5(a-1).
\]
The subsequent return traverses the same positions in the opposite direction.
Over the eighteen activations,
\begin{equation}
 \sum N=90a,\qquad\sum r=90,\qquad
 \sum q=10(a-1)=
 \begin{cases}30,&a=4,\\40,&a=5.\end{cases}
 \label{bim09:ocupacion-30-40}
\end{equation}
The accumulated integer product is, respectively, $360$ or $450$.
The net winding is zero; the sum of quotients and the eighteen
incidences preserve the history of the reading. Each of the three
signature classes contains representatives of both transverse values.
The same signature and the same transport class of that finite quotient
can therefore coexist with distinct arithmetic occupations.

These results connect preserved memory with concrete
operations: permutation of half-cycles, boundary differences,
arithmetic quotients, and incidence records. The physical-mathematical
identification of any subsequent reader preserves its own
maps; the identities proved here belong to the internal dynamics
of the product channel.

% Brief insertion fragment. Development and provenance in the same folder.
\subsubsection{Arithmetic memory of the multiplicative channel}
\label{prd-arit:seccion}

Preservation of the regional signature extends to the full
arithmetic evaluation of APP. For an oriented position
\(x=(i,j,d)\in\{1,\ldots,9\}^2\times\{N,E,S,O\}\), let
\(\chi(x)=(a,b,\epsilon)\), where \(a\) is the fixed coordinate,
\(b\) the moving coordinate, and \(\epsilon\) its orientation:
\[
 \chi(i,j,N)=(j,i,-1),\quad \chi(i,j,S)=(j,i,+1),\qquad
 \chi(i,j,E)=(i,j,+1),\quad \chi(i,j,O)=(i,j,-1).
\]
The axis label \(\iota\in\{i,j\}\) preserves which of the
two original coordinates moves. Writing
\(\langle b\rangle_9=1+((b-1)\bmod9)\), the operators and
evaluation are
\begin{align}
 \overline P(a,b,\epsilon)&=(a,\langle b+\epsilon\rangle_9,\epsilon),
 &\overline H(a,b,\epsilon)&=(a,b,-\epsilon),\label{prd-arit:PH}\\
 n(a,b)&=ab,&r(a,b)&=\langle ab\rangle_9,\qquad
 q(a,b)=\frac{ab-r(a,b)}9.\label{prd-arit:rq}
\end{align}

\begin{proposition}[Exact representation of arithmetic transport]
\label{prd-arit:minimal}
The identities \(\chi P=\overline P\chi\) and
\(\chi H=\overline H\chi\) hold. The representation \(\chi\)
has \(162\) states and is the coarsest of the representations
stable under \(P,H\) that preserve the integer product or,
equivalently, the pair \((r,q)\). Additionally preserving \(\iota\)
recovers the \(324\) original oriented positions.
\end{proposition}
\begin{proof}
Advance modifies only \(b\), and the half-turn only
\(\epsilon\), proving the two transport identities.
The nine observations \(n_k=n(\overline P^k\chi(x))\),
\(0\leq k<9\), run through \(a,2a,\ldots,9a\). Thus one recovers
\(a=\min_k n_k\), \(b=n_0/a\), and the sign through
\((n_1/a-n_0/a)\bmod9\in\{1,8\}\). Every stable equivalence
that preserves \(n\) preserves these nine observations and therefore
distinguishes all states \((a,b,\epsilon)\).
The identity \(n=r+9q\) proves equivalence of the two
evaluations. Each fiber of \(\chi\) contains the two transposed
routes: \((i,j,d)\mapsto(j,i,d^\top)\), with
\(N^\top=O,O^\top=N,E^\top=S,S^\top=E\). The axis label distinguishes
these two preimages and allows position and direction to be reconstructed.
\end{proof}

The TRIT selector activates the product at \(t\equiv2,5,8\pmod9\),
before each advance. The half-turn acts at the end of every
\(27\) instants. If \(y_{t-1}\) denotes the preceding arithmetic
state, the increment and its accumulated memory are determined by
\begin{equation}
 \eta_q(t)=\mathbf1_{\{t\equiv2\ (3)\}}q(y_{t-1}),\qquad
 \lambda_q(t)=\lambda_q(t-1)+\eta_q(t),\quad \lambda_q(0)=0.
 \label{prd-arit:incremento}
\end{equation}
The bound \(0\leq q(a,b)\leq a-1\) holds. To invert an
update, one first undoes the corresponding half-turn,
then the advance, and finally subtracts the quotient recovered from the
preceding state. The record retains phase, axis, event order, and
memory; \(\lambda_q\) is an additive cocycle of this transport.

Each block of \(27\) contains nine advances and traverses the
moving coordinate once. Hence the return of position and orientation
at \(54\), and of phase at \(108\), coexist with
\begin{equation}
 \lambda_q(108N)=NQ_\times^+(a),\qquad
 Q_\times^+(a)=\frac{180a-4\sum_{b=1}^9r(a,b)}9,
 \qquad N\geq0.
 \label{prd-arit:retorno}
\end{equation}
Indeed, each block sums the products \(a,2a,\ldots,9a\);
summing \(n=r+9q\) over the four blocks gives the formula. If
\(a\) is coprime to \(9\), the residues sum to \(45\), and
\(Q_\times^+(a)=20(a-1)\). In window \(j\), counted from
\(j=0\), the relative calendar sign is
\(\sigma_j=(-1)^{\lfloor j/3\rfloor}\), and the effective
orientation is \(\epsilon_j=\epsilon_0\sigma_j\).
The charge signed by \(\sigma_j\) vanishes; the charge weighted
by the effective orientation likewise vanishes after multiplication by
\(\epsilon_0\). The ledger preserves the initial orientation, the accumulated
charge, and its increments. These quantities belong
to the multiplicative cursor; the readers of a trajectory that
alternates between both sheets act on the record of that trajectory.


The finite emission, its coupling of signatures, and ternary reduction have
successive types:
\[
 \mathcal S^2\xrightarrow{s\times e}
 \operatorname{im}s\times\operatorname{im}e
 \xrightarrow{G}\mathcal U_6
 \xrightarrow{\Psi_6}W_6\hookrightarrow\mathbb F_3^6.
\]
Section~\ref{sec:extension} proves the factorization of the
\(104\,976\) seeds into \(18\cdot26=468\) emissions and the
\(243\) visible words in the ambient space of \(729\) elements. The
maps \(G\) and \(\Psi_6\) are those in
\eqref{eq:acoplamiento} and~\eqref{eq:psi-catalogo}; the record preserves
the arithmetic preimages together with this finite image.

\subsubsection{Compatibility of quotients and refinement cylinders}

Prolongation compares states that simultaneously preserve remainders,
quotients, and boundaries. In the six-coordinate linear chart of the record,
the matrix
\[
 A_H=\begin{pmatrix}
 2&2&2&1&2&1\\2&1&2&2&1&1\\1&0&0&1&0&2\\
 0&0&1&1&1&0\\2&2&2&0&2&1\\2&1&2&1&0&0
 \end{pmatrix},\qquad\det A_H=1,
\]
gives the division map
\begin{equation}
 \mathcal H:\mathbb Z^6\longrightarrow
 \{0,1,2\}^6\times\mathbb Z^6,\qquad
 x\longmapsto(u,x^+),\quad
 y=xA_H,\quad u=[y]_3,\quad x^+=(y-u)/3.
 \label{ii:tpk-hensel-local}
\end{equation}
Rows are interpreted as vectors. The pair \((u,x^+)\) preserves
\(y=u+3x^+\) and permits recovery of
\(x=(u+3x^+)A_H^{-1}\). The ternary orientation of \(u\) is then obtained
through \(\rho_3^{\rm or}\), preserving the preceding quotient.
This makes the reversible content of the local chart explicit.

The cylindrical chart relates a ternary prefix of length \(L\), with
integer index \(N\), to \(K\) decimal triads, with index \(D\).
The corresponding half-open cylinders have rational endpoints
\[
 I_3(N,L)=\left[\frac N{3^L},\frac{N+1}{3^L}\right),\qquad
 I_{1000}(D,K)=\left[\frac D{1000^K},\frac{D+1}{1000^K}\right).
\]
Define the boundary integers
\[
 a=N1000^K-D3^L,\qquad b=(N+1)1000^K-D3^L.
\]

\begin{lemma}[Integer update of compatible cylinders]
\label{ii:tpk-cilindros-enteros}
The preceding cylinders intersect if and only if
\(a<3^L\) and \(b>0\). Adjoining a ternary block
\(u\in\{0,\ldots,728\}\) gives \(N'=729N+u\) and
\(L'=L+6\). For \(\Delta K\in\{0,1\}\), the boundaries are
updated by
\[
 \begin{array}{ll}
 \Delta K=0:&a'=729a+u1000^K,\\[2pt]
 \Delta K=1:&D'=1000D+\delta,\quad
 a'=729000a+u1000^{K+1}-\delta\,729\,3^L,
 \end{array}
\]
where \(\delta\in\{0,\ldots,999\}\) is the triad of the
compatible prolongation. In both cases,
\(b'=a'+1000^{K+\Delta K}\).
\end{lemma}
\begin{proof}
Two half-open intervals intersect exactly when each left
endpoint is less than the other's right endpoint. Multiplying those two
inequalities by \(3^L1000^K\) gives \(a<3^L\), \(b>0\).
Substituting \(N'\), \(L'\), and \(D'\) into
\(a'=N'1000^{K+\Delta K}-D'3^{L'}\) produces the two recurrences.
Subtracting the definitions of \(b'\) and \(a'\) gives the last identity.
\end{proof}

These operations act on integer prefixes and their compatible
restrictions. The Archimedean reading appears when the resulting
cylinder system is evaluated. The prolongation domain also incorporates
orientation, carry, and the joint boundary signature of
\eqref{ii:tpk-fibra}; interval comparison is one component of
this joint compatibility.

\subsubsection{Extension relations and joint nonadic connection}

For an observable route \(r:q\to q'\), prolongation has type
\[
 \operatorname{Ext}_r\subseteq\mathcal F_q\times\mathcal F_{q'},
 \qquad
 \operatorname{Ext}_{r_2}\circ\operatorname{Ext}_{r_1}
 \subseteq\operatorname{Ext}_{r_2r_1}.
\]
The identities satisfy
\(\Delta_{\mathcal F_q}\subseteq
\operatorname{Ext}_{\operatorname{id}_q}\).
It preserves the truncations, sheets, orientations, carries, and
boundary incidences of the two histories. The identities and
admissible concatenations form the category of TPK states and prolongations.
This relational structure admits local fibers with more than one
continuation and allows their persistence to be imposed jointly.

Let \(\mathcal H_m\) be the set of admissible histories of depth
\(m\), with truncations \(p_{n,m}\). Survival is expressed by
\[
 \operatorname{Surv}_m^{(N)}=p_{N,m}(\mathcal H_N),\qquad
 \operatorname{Surv}_m=\bigcap_{N\geq m}
 \operatorname{Surv}_m^{(N)},\qquad
 X_\chi=\varprojlim_m\operatorname{Surv}_m^{(\chi)}.
\]
The compatible section simultaneously preserves all its prefixes. The
local connection relations \(\mathcal R_k\) compose into
\begin{equation}
 \Gamma_{9,k}=\mathcal R_{k+8}\circ\cdots\circ\mathcal R_k,\qquad
 \Gamma_{9,k}:\widehat X_k\dashrightarrow\widehat X_{k+9}.
 \label{ii:tpk-holonomia}
\end{equation}
The indicated domain is that of compatible histories. On an individualized
history, the nine relations represent its effective transport;
the system retains all admissible prolongations before this
individualization.

The return law is written
\[
 g(\Gamma_{9,k}x)=g(x),\qquad
 q_{\rm mem}(\Gamma_{9,k}x)=q_{\rm mem}(x)+1.
\]
The phase returns and the history is prolonged. The prefixes
\(w_6,w_{12},w_{18},w_{24},w_{30}\), the boundary \(R_{36}\), and the
charts \(G_9\) are states, extensions, or presentations of this
construction, with types distinct from the holonomy \(\Gamma_9\).
Their regional realization is developed in section~\ref{sec:generacion};
the composition of turns and their representations are proved in
\ref{mono-ampl:def-vuelta}--\ref{mono-ampl:representacion}.

\subsubsection{Phase incidence and arithmetic origin of the channels}

The selector \(M_{\rm ph}\) induces the classes
\[
 H_+=\{1,4,7\},\quad H_-=\{2,5,8\},\quad
 H_0=\{3,6,9\},\quad H_{\rm act}=H_+\sqcup H_-.
\]
The phase origin fixes the terminal representative
\(\star_{\rm ph}=9\) and the set
\(O_{\rm ph}=D_9\setminus\{\star_{\rm ph}\}\).
Write \(P_2(H_{\rm act})\) for the unordered pairs of active
phases and define
\[
 \mathcal F_6^{\rm ph}=H_{\rm act}\times P_2(H_{\rm act}),\qquad
 \mathcal F_8^{\rm ph}=O_{\rm ph}\times P_2(H_{\rm act}),\qquad
 \Theta_{\rm ph}=D_9\times P_2(H_{\rm act}).
\]

\begin{theorem}[Incidence of the phase semicorona]
\label{ii:tpk-semicorona}
At the declared window origin,
\[
 |\mathcal F_6^{\rm ph}|=90,\qquad
 |\mathcal F_8^{\rm ph}|=120,\qquad
 |\Theta_{\rm ph}|=135.
\]
The oriented boundary, pointed return, and lateral closure,
\[
 \partial_{\rm ph}=\Theta_{\rm ph}\times\{-1,+1\},\quad
 E_{\rm ph}=\partial_{\rm ph}\sqcup\{\infty\},\quad
 L_{\rm ph}=H_0\times P_2(H_{\rm act}),\quad
 P_{\rm ph}=\partial_{\rm ph}\sqcup L_{\rm ph},
\]
satisfy
\[
 |\partial_{\rm ph}|=270,\quad |E_{\rm ph}|=271,\quad
 |L_{\rm ph}|=45,\quad |P_{\rm ph}|=315.
\]
Translation of the origin transports all these sets by bijections
and preserves their incidences.
\end{theorem}
\begin{proof}
There are six active phases and \(\binom62=15\) unordered pairs.
Products by \(6,8,9\) give \(90,120,135\). Orientation
doubles the last set, the pointed return adds one element, and the
three neutral phases contribute \(3\cdot15=45\) lateral incidences.
The \(C_9\) action simultaneously transports the selector, the
terminal representative, and each factor of the Cartesian products.
\end{proof}

The difference of active cardinalities is \(120-90=30\). The normalized
defect therefore has the exact expression
\begin{equation}
 \frac{30}{120}-\frac{30}{90}=-\frac1{12}.
 \label{ii:tpk-defecto-incidencia}
\end{equation}
Here \(90\) and \(120\) are cardinalities of incidences produced by the
phase rule. The subsequent hexadic--octadic realization transports this
incidence to its combinatorial objects. The angular coordinates of these
objects are established through their realization maps.

\subsubsection{Temporal record, signed extraction, and reversible coordinates}

Let \(E_m=\{9(m-1)+1,\ldots,9m\}\). Restricting a history
to twelve windows preserves in each its subroute \(\tau_m\),
orientation \(\sigma_m\), boundary carry \(\kappa_m\), and local
record \(\Lambda_m\):
\[
 \mathcal R_{12}(x)
 =\bigl((E_m,\tau_m,\sigma_m,\kappa_m,\Lambda_m)\bigr)_{m=1}^{12}.
\]
The domain of \(\mathcal R_{12}\) is histories; the elements of
the group \(C_{12}\) are sector classes, and
\(\Theta_{108}\) preserves their calendar. This distinction keeps
the data preserved by each realization visible.

The signed extraction of section~\ref{sec:k-registro} factors through
\begin{equation}
 \mathcal X_{\rm TPK}^{\rm term,mem}
 \xrightarrow{\mathcal R_{12}}\mathcal R_{12}^{\rm mem}
 \xrightarrow{\mathcal E_{108}^{90,120}}
 \mathbb Z^{12}\times\mathbb Z^{12}\times\mathbb Z
 \xrightarrow{\Pi_H}\mathbb Z^{12}.
 \label{ii:tpk-extraccion-tipada}
\end{equation}
The second map extracts the channels
\((b^{90},b^{120},Q_{\rm TPK})\); the third constructs the signed
vector \(U\). The rule \(\varepsilon_9=M_{\rm ph}\) belongs to
phase selection, whereas \(\mathcal E_{108}^{90,120}\) acts on
the full channel record. Their domains and outputs distinguish them.

Equations~\eqref{eq:k-sumas-armonicas} and
\eqref{eq:k-bloques-firmados} give the arithmetic action of \(\Pi_H\).
The transformation \(\mathcal H_{12}\), constructed from the
three step-three orbits, then reversibly relates \(U\) to the
record \(K\). The cyclic differences and total charge provide
subsequent readings of the same record. This preserves the common
provenance of the regional coordinates and the signed coordinate before the
integer closure determining \(\alpha\).

\subsubsection{Reduced transfer and memory preserved by the dynamics}

Once the emission factorization has been constructed,
\(\mathcal U_6\simeq\mathcal A_+\times\mathcal A_\times\), with
\(|\mathcal A_+|=18\) and \(|\mathcal A_\times|=26\), there are
two projections on functions of the catalog:
\[
 (E_+f)(s,e)=\frac1{18}\sum_{s'\in\mathcal A_+}f(s',e),\qquad
 (E_\times f)(s,e)=\frac1{26}\sum_{e'\in\mathcal A_\times}f(s,e').
\]
The signature alphabets \(\mathcal A_+,\mathcal A_\times\) are
distinct from the seed sets of the two cursors.
Each projection averages the active coordinate and preserves the other. With
\(P_{+1}=E_+\), \(P_{-1}=E_\times\), \(P_0=I\), the transfer
kernel on \(\mathcal U_6\times C_9\) is
\[
 \mathcal K_{\rm ph}((u,r),(v,r+1))=P_{M_{\rm ph}(r)}(u,v).
\]
The two projections commute and are idempotent. Thus their composition
during one phase turn satisfies
\begin{equation}
 \mathcal K_{\rm ph}^{9}\big|_r=E_+E_\times,\qquad
 (E_+E_\times)(u,v)=\frac1{468}.
 \label{ii:tpk-renovacion}
\end{equation}
This operator is a reduced-memory representation on emissions already
constructed. The holonomy \(\Gamma_9\) also preserves the
coordinates of the history on which that representation is obtained.

In an isometric realization \(U_{\Gamma_9}\) of transport, and for
an isometric inclusion \(J\) of the observed coordinate, its two
components are
\[
 T=J^*U_{\Gamma_9}J,\qquad
 \eta=(I-JJ^*)U_{\Gamma_9}J.
\]
One obtains directly
\[
 U_{\Gamma_9}J=JT+\eta,\qquad J^*\eta=0,\qquad
 I-T^*T=\eta^*\eta.
\]
The component \(\eta\) is computed from the same realized update
and the declared projection. The telescoping identity
\eqref{eq:telescopia} preserves both the contributions expressed in the
orthogonal complement and the terminal term. This separation is the
operator basis for the subsequent distinction between recorded information,
entropy of an observation, and dimensional realization of action.
